\documentclass[11pt]{article}
\usepackage[margin=0.85in]{geometry}
\usepackage{amsmath,amssymb,booktabs,microtype}
\usepackage[T1]{fontenc}
\usepackage{lmodern}
\usepackage[hidelinks]{hyperref}
\setlength{\parindent}{0pt}
\setlength{\parskip}{6pt}
\newcommand{\E}{\mathbb E}
\begin{document}
\begin{center}
{\Large Housing over three stages of life}\\[4pt]
{\large A test of the proposed change}\\[3pt]
September 9, 2026
\end{center}

Adding a mature working period can generate the intended housing comparison
with an 80\% mortgage. It also permits analytical household solutions. It
does not, by itself, deliver a simple formula for equilibrium prices or a
fertility result along a policy transition. This calculation tests the change;
it does not replace the model in the seminar deck.

\textbf{1. What changes.}
A household is young, mature, then old. Its incomes are $y^y,y^m,y^o$,
and its initial financial endowment is $b$; write $w=y^y+b$.
Income and wealth may differ across households. Fertility $n$ is chosen when
young. Children require goods and space in both family stages:
\begin{align*}
u^y&=\log(c^y-\chi n)+\alpha\log(h^y-\kappa n)+\vartheta\log n,\\
u^m&=\log(c^m-\chi n)+\alpha\log(h^m-\kappa n),\\
u^o&=\log c^o+\gamma\log h^o+\omega_B\log e.
\end{align*}
Lifetime utility is $u^y+\beta u^m+\beta^2u^o$, plus the existing tenure
taste. Here $e$ is the estate, including the retained home's liquidation value.
Tenure is retained, homes can be resized, and each home's size is bounded by
$H_R$ or $H_O$, with $H_R<H_O$.

At a stationary price $P$, let $q$ be the price of a bond paying one unit next
stage, $\tau$ the property-tax rate, and $T$ the common rebate. The current
cost of housing services is $p=DP$, where $D=1-q+q\tau$.
Net financial assets entering maturity and old age are $a_m$ and $a_o$.
An owner's budgets are
\begin{align*}
c^y+qa_m+(1+q\tau)Ph^y&=w+T,
&qa_m+\phi Ph^y&\ge0,\\
c^m+qa_o+(1+q\tau)Ph^m&=a_m+Ph^y+y^m+T,
&a_o&\ge0,\\
c^o+qe+ph^o&=a_o+Ph^m+y^o+T,
&e&\ge Ph^o.
\end{align*}
A renter pays $ph$ at each age, has no housing-sale receipts, and faces
$a_m,a_o\ge0$. Its old budget is $c^o+qe+ph^o=a_o+y^o+T$.
The original logistic tenure choice and the common rebate are retained.

\textbf{The economic difference is the repayment date.}
When the young mortgage binds, resources on entry to maturity are
\[
\boxed{\quad y^m+T+(1-\phi/q)Ph^y.\quad}
\]
The entire mortgage is repaid, including interest. Mature earnings can cover
that repayment before retirement. The tested restriction $a_o\ge0$ forbids
negative net financial wealth at retirement; it is an additional institutional
choice. Where the derived $a_o$ is strictly positive, permitting mature
refinancing leaves the household optimum unchanged.

\newpage
\textbf{2. How much can be solved explicitly?}

Consider an owner whose young mortgage binds, who saves when mature, and whose
housing caps and old estate floor are slack. Let $x=c^y-\chi n$ and
$s=h^y-\kappa n$ be young consumption and space after child needs. Define
\begin{gather*}
W=w+T,\qquad U=y^m+qy^o+(1+q)T,\\
L=(1-\phi+q\tau)P,\quad \delta=(1-\phi/q)P,\quad
t=\chi+\kappa p,\quad t_L=\chi+\kappa L,\quad d=t-\kappa\delta,\\
K=1+\gamma+\omega_B,\qquad G=1+\alpha+\beta K.
\end{gather*}
Let $x^m=c^m-\chi n$ and $s^m=h^m-\kappa n$ denote mature adult consumption
and space. Mature resources after child needs are $B=U+\delta s-dn$.
The continuation choices are
\[
x^m=\frac BG,\quad s^m=\frac{\alpha x^m}{p},\quad
c^o=\frac{\beta x^m}{q},\quad h^o=\frac{\gamma\beta x^m}{qp},\quad
e=\frac{\omega_B\beta x^m}{q^2}.
\]

\textit{Fixed fertility.} Set $A_y=W-t_Ln$ and $A_m=U-dn$.
Young space solves a strictly concave, one-variable problem:
\[
\max_s\ \log(A_y-Ls)+\alpha\log s+\beta G\log(A_m+\delta s).
\]
Its first-order condition is the quadratic
\[
\boxed{\begin{aligned}
0={}&\alpha A_yA_m+
\big[\delta A_y(\alpha+\beta G)-(1+\alpha)LA_m\big]s\\
&-\delta L(1+\alpha+\beta G)s^2.
\end{aligned}}
\]
Select the root with $s>0$, $A_y-Ls>0$, and $A_m+\delta s>0$.
This holds at general $\phi$, including $\phi=.8$.

\textit{Chosen fertility.} Put $r=\beta x/x^m$ and $E=1+\alpha+\vartheta$.
Then
\[
s=\frac{\alpha x}{L-\delta r},\qquad
n=\frac{\vartheta x}{t_L+dr},\qquad
x=\frac{W+Ur}{E+\beta G},
\]
where $r$ solves
\[
\left(1+\frac UW r\right)
\left(1+\frac{\alpha L}{L-\delta r}
+\frac{\vartheta t_L}{t_L+dr}\right)=E+\beta G.
\]
Clearing denominators gives a cubic. Its economically relevant root is unique.
The income ratio $U/W$ is strictly decreasing in that root: higher mature
resources relative to young cash lower $r$. Thus a condition $r<r_*$ can be
written as an explicit income-ratio inequality without solving the cubic.

These formulas must satisfy $0<r<q$, the housing limits, and
\[
a_o=\left(\frac{\beta K}{q}-\frac\alpha D\right)x^m
-(y^o+T)-\kappa Pn>0,\qquad \omega_BD>q\gamma.
\]
In particular, $\beta KD>\alpha q$ is necessary for this saving regime.
Higher mature earnings alone cannot repair a failure of that inequality.

\newpage
\textbf{3. The planner and fertility comparisons remain short.}

At the stationary reference, match the common distribution of types and tenure
across the three current ages. Bars denote means under this distribution.
Hold fertility, future real opportunities and old estates fixed. The planner
chooses all current consumption and housing, giving equal weight to the
remaining utility of each current household. It respects both current resource
constraints and each household's tenure cap, and may relax private finance limits.

After subtracting the fixed child needs, total current goods are
$\bar x^y+\bar x^m+\bar c^o$. Hence every household receives adult consumption
\[
X=\frac{\bar x^y+\bar x^m+\bar c^o}{3}.
\]
For a type with fertility $n_i$ and cap $H_i$, planner housing is
\[
h_i^{y,F}=h_i^{m,F}=\min\{H_i,\kappa n_i+z\},\qquad
h_i^{o,F}=\min\{H_i,(\gamma/\alpha)z\},
\]
where $z$ clears the housing stock. Title and financial-claim adjustments keep
future entering resources and net estates fixed. The required current transfers
sum to zero because current goods and housing totals are unchanged.

When planner caps are slack, the young receive more housing exactly when
\[
\boxed{\quad
\alpha(\bar s^m+\bar h^o)>(\alpha+\gamma)\bar s^y.
\quad}
\]
At $\alpha=\gamma$, mature and old adult space together must exceed twice
young adult space. The mature cohort is part of the comparison.

A particularly simple sufficient statement also allows binding caps. If
$\alpha\ge\gamma$, fertility is positive and some total capacity is unused,
the planner gives each family age more than one third of the housing stock.
Therefore, if the competitive young hold at most one third, their aggregate
housing rises. This is a utilitarian allocation comparison, not a Pareto claim.

\textit{Private fertility after a current allocation change.}
Children still require goods and housing when mature. The private condition is
\[
\frac{\vartheta}{n_i}
=\frac\chi{x_i^y}+\frac{\alpha\kappa}{s_i^y}
+\beta\left(\frac\chi{x_i^m}+\frac{\alpha\kappa}{s_i^m}\right).
\]
Give a parent its new current consumption and housing bundle, while preserving
its continuation problem. At the original fertility, the future marginal cost
is the same in both comparisons. Strict concavity then gives
\[
\boxed{\quad n_i^{\rm new}>n_i
\quad\Longleftrightarrow\quad
\chi\left(\frac1{x_i^y}-\frac1X\right)
+\alpha\kappa\left(\frac1{s_i^y}-\frac1{s_i^{y,F}}\right)>0.\quad}
\]
Thus more housing and weakly more adult consumption increase fertility.
A housing gain alone need not do so. This is a household comparison at the
same continuation opportunities; a market policy changes prices and future
choices as well, and needs its own argument.

\newpage
\textbf{4. What the equilibrium check establishes.}

There is an analytical stationary-equilibrium construction with arbitrary
fixed $\phi\in(0,1)$, including $.8$. Young and old endowments have bounded
heterogeneous support. Mature earnings are sufficiently high relative to
young cash resources. Take $\alpha\ge\gamma$. At zero tax, also require
\[
\omega_B(1-q)>q\gamma,\qquad
\beta(1+\gamma+\omega_B)(1-q)>q\alpha.
\]
Finite inequalities bound the required mature earnings, rental capacity and
owner capacity. On an explicit positive price interval they imply: young
owners borrow maximally; mature owners save; both later owner homes exceed
the young home; and mature and old renters occupy their rental cap while
young renters remain below it. The cap still affects young fertility through
the family's future mature housing needs.

Both tenure-specific fertility choices exceed replacement at the lower price
and fall below it at the upper price. These inequalities hold type by type,
so endogenous logistic tenure shares preserve them after aggregation.
Continuity gives a stationary equilibrium inside that interval. At such an
equilibrium the young hold less than one third of housing. The fixed-fertility
planner described on page 3 gives them more than one third.

This is an analytical existence proof with finite primitive bounds, not a
numerical equilibrium followed by a neighbourhood argument. It is nevertheless
a deliberately strong sufficient construction: the earnings bounds can be
large, and it does not establish the ordering in every possible equilibrium
outside the constructed price interval.

Further finite bounds on mature earnings and rental capacity make every young
household receive both more consumption and more housing under the dated planner.
The private fertility test on page 3 is then strictly positive for every young
parent. This establishes a conditional fertility response at unchanged
continuation opportunities, not a future market-clearing allocation.

\textit{Taxes.} The strict construction extends to sufficiently small positive
rebated property taxes. That extension uses continuity and a two-variable
fixed-point argument; the maximal admissible tax rate has not been derived.
Independently, every stationary equilibrium obeys the explicit rebate bound
\[
T\le
\frac{\tau(\bar w+q\bar y^m+q^2\bar y^o)}
{(1-q)[3q-\tau(1+2q)]},
\qquad \tau<\frac{3q}{1+2q}.
\]
This bound alone does not establish the desired household regime at a chosen tax.

\textit{A less extreme household check.} With $q=.5$, $\tau=.1$, $\phi=.8$,
$\alpha=\gamma=\vartheta=\kappa=1$, $\omega_B=2$, and
$\beta\in[.49,.5]$, an explicit conditional income interval delivers binding
young credit, mature saving and larger old homes with mature resources between
$2.5$ and $2.8$ times young resources. This check also requires
$\chi/(\kappa p)\ge10$ and comparable young and old cash resources.
It verifies household feasibility; it is not the stationary construction above.

\textbf{Assessment.} The extra earning stage can generate the intended housing
mechanism with conventional origination finance. Fixed-fertility household
choices are quadratic; endogenous fertility is cubic. With heterogeneous types
and tenure choice, stationary prices still solve aggregate clearing equations,
and have no demonstrated closed-form solution. A short general primitive
condition covering housing, fertility and a policy transition has not been
obtained. The model change therefore merits consideration, but this calculation
does not justify replacing the seminar theory wholesale.

\small Full budgets, derivative checks and the finite equilibrium inequalities
are retained together in the existing simplified-OLG evidence folder,
in \nolinkurl{three_stage_tractability_evidence.md}. No numerical model runs were used.
\end{document}
